% TOP_PROBLEM: {"problemId": "problem.p-versus-nc-1", "canonicalTitle": "P versus NC^1", "releaseRank": 69, "primaryDomain": "theoretical_computer_science", "primaryDomainLabel": "Theoretical computer science", "displayStatus": "open", "releaseStatus": "open"}
% !TeX program = lualatex
% Definition and sources only; this file is not an LLM solution attempt.
\documentclass[11pt]{article}
\usepackage[margin=25mm]{geometry}
\usepackage{fontspec}
\setmainfont{Latin Modern Roman}
\usepackage{amsmath,amssymb,mathtools}
\usepackage{unicode-math}
\setmathfont{Latin Modern Math}
\usepackage{xurl,hyperref}
\hypersetup{colorlinks=true,urlcolor=blue}
\setcounter{secnumdepth}{0}
\setlength{\parindent}{0pt}
\setlength{\parskip}{5pt}
\setlength{\emergencystretch}{3em}
\begin{document}
\section*{\texorpdfstring{P versus NC\textasciicircum{}1}{P versus NC\textasciicircum{}1}}\label{69}
\subsection{Definitions and mathematical statement}
A De Morgan formula is a Boolean circuit whose computation graph is a tree, using AND, OR, and NOT. Its size counts its gates or symbols, up to the usual polynomially equivalent conventions. The target is a uniformly polynomial-time computable family $f_n:\{0,1\}^n\to\{0,1\}$ with
\[
 \operatorname{FormulaSize}(f_n)\text{ not bounded by any polynomial in }n.
\]
The formulas themselves are nonuniform. This is the assertion $\mathsf P\nsubseteq\mathsf{NC}^1$ for nonuniform $\mathsf{NC}^1$; an absence of a uniform shallow-circuit construction would be weaker.
\par\smallskip\textit{Formulation and concept references:} \href{https://eccc.weizmann.ac.il/report/2013/190/revision/11/download}{[S1]}.

\subsection{Short English statement}
Is there an efficiently computable Boolean function family that no polynomial-size formula family can represent?
\subsection{Sources}
\begin{itemize}
\item[S1]Gavinsky, Meir, Weinstein, and Wigderson, ECCC TR13-190, revision 11, 2017.
\url{https://eccc.weizmann.ac.il/report/2013/190/revision/11/download}.
\textit{Formulation source.}
\item[S2]Toward Better Depth Lower Bounds: Strong Composition of XOR and a Random Function.
\url{https://arxiv.org/abs/2410.10189}.
\textit{Further reference; not a proof endorsement.}
\item[S3]Toward Better Formula Lower Bounds: The Composition of a Function and a Universal Relation.
\url{https://epubs.siam.org/doi/10.1137/15M1018319}.
\textit{Further reference; not a proof endorsement.}
\end{itemize}
\end{document}
